% Standalone circuitikz version of esp32-multimeter_schem, applying the
% lessons from ../potentiometer_multimeter_tikz/ and
% ../led_resistor_esp32_tikz/ (see ../../circuit-diagrams-tikz.md).
% Build with build.sh in this directory.
%
% Same hand-drawn ESP32 box as led_resistor_esp32.tex (gotcha #7: plain
% shapes, not circuitikz IC/meter bipoles -- a voltmeter bipole here would
% draw an unwanted needle same as the ohmmeter did). No wire color (the
% schemdraw version tried colored wires and was asked to drop it in favor
% of plain text labels naming the probe and port).
\documentclass[tikz,border=4pt]{standalone}
\usepackage{circuitikz}
\usepackage[scaled]{helvet}
\renewcommand{\familydefault}{\sfdefault}

\definecolor{navy}{HTML}{2C3E50}

% Gotcha #8: match any bipole body thickness to plain wires.
\ctikzset{bipoles/thickness=1}

\begin{document}
\begin{circuitikz}[american, color=navy, thick, line width=1pt]
    \large

    % ESP32 box with two pins on the left side
    \draw (0,0) rectangle (2.4,6);
    \node at (1.2,3) {\shortstack{ESP32\\FireBeetle}};
    \coordinate (A4PIN) at (0,5);
    \coordinate (GNDPIN) at (0,1);
    \draw (A4PIN) -- ++(-0.6,0) coordinate (A4TIP);
    \draw (GNDPIN) -- ++(-0.6,0) coordinate (GNDTIP);
    \node[right] at (0.1,5) {A4};
    \node[right] at (0.1,1) {GND};

    % multimeter, in series between the two pins -- same x for both arms
    \coordinate (CHAINX) at ($(A4TIP)-(2,0)$);
    \coordinate (TOP) at (CHAINX |- A4TIP);
    \coordinate (BOT) at (CHAINX |- GNDTIP);

    \node[draw, circle, minimum size=1.3cm, line width=1pt] (METER)
        at ($(TOP)!0.5!(BOT)$) {$V$};

    \draw (A4TIP) -- (TOP) node[circ]{} -- (METER.north);
    \draw (GNDTIP) -- (BOT) node[circ]{} -- (METER.south);

    % each arm labeled with its probe and the multimeter port it plugs
    % into -- offset away from the arm itself (gotcha #12: anchors sit
    % flush against the drawn geometry, not with any built-in clearance)
    \node[above] at ($(A4TIP)!0.5!(TOP)+(0,0.15)$) {\shortstack{red probe\\V$\Omega$mA port}};
    \node[below] at ($(GNDTIP)!0.5!(BOT)+(0,-0.15)$) {\shortstack{black probe\\COM port}};

\end{circuitikz}
\end{document}
